\section*{Chapter \ref{ch:b3:seismology-imaging} --- Seismology and Earth Imaging}

\begin{solution}{exo:b3:seismology-imaging:1}
(a) $\sin i_0 = pv_0/R = 500\times5.80/6371 = 0.4552$, so $i_0 = \ang{27.1}$.
(b) At \qty{771}{km}, $r = \qty{5600}{km}$ and $r/v = 5600/11.066 =
\qty{506.1}{s.rad^{-1}}$; at \qty{871}{km}, $r = \qty{5500}{km}$ and $r/v =
5500/11.245 = \qty{489.1}{s.rad^{-1}}$. The ray turns where $r/v = p = 500$:
$z_t = 771 + 100\times(506.1 - 500)/(506.1 - 489.1) = 771 + 36 =
\qty{807}{km}$.
\end{solution}

\begin{solution}{exo:b3:seismology-imaging:2}
(a) $(455.5 - 369.6)/10 = \qty{8.59}{s}$ per degree, and
$8.59\times180/\pi = \qty{492.2}{s.rad^{-1}}$. (b) One degree of arc is
$6371\pi/180 = \qty{111.19}{km}$, so the wave front sweeps the surface at
$111.19/8.59 = \qty{12.94}{km.s^{-1}}$, far faster than any rock at the
surface: the ray arrives steeply. (c) At the turning point $v = r_t/p =
5450/492.2 = \qty{11.07}{km.s^{-1}}$, a lower-mantle P speed.
\end{solution}

\begin{solution}{exo:b3:seismology-imaging:3}
(a) $\mu_1 = \rho_1\beta_1^2 = 2700\times3500^2 = \qty{3.31e10}{Pa}$
(\qty{33.1}{GPa}); $\mu_2 = 3300\times4500^2 = \qty{6.68e10}{Pa}$
(\qty{66.8}{GPa}). (b) $\lambda = cT = 3.79\times20 = \qty{75.8}{km}$, 2.17
times the layer thickness: the wave reaches well into the half-space, which is
why its speed is already above $\beta_1$. (c) $4000/3.384 = \qty{1182}{s}$ and
$4000/4.333 = \qty{923}{s}$: the 20-s waves arrive \qty{259}{s}, or
\qty{4.3}{min}, after the 100-s waves.
\end{solution}

\begin{solution}{exo:b3:seismology-imaging:4}
(a) ${}_0S_2$: $1/(53.89\times60) = \qty{3.093e-4}{Hz} = \qty{0.3093}{mHz}$;
${}_0T_2$: $1/(44.01\times60) = \qty{0.3787}{mHz}$. (b) $1440/53.89 = 26.7$
and $1440/44.01 = 32.7$ oscillations a day. (c) The difference is
$\qty{0.0694}{mHz}$, whose inverse is \qty{14400}{s}: at least \qty{4.0}{h} of
record, in practice much more, since the peaks are broadened by attenuation and
split.
\end{solution}

\begin{solution}{exo:b3:seismology-imaging:5}
(a) The chord between two points of the sphere separated by the angle $\Delta$
at the centre has the length $2R\sin(\Delta/2)$, travelled at $v$.
(b) At \ang{30}: $2\times6371\times\sin\ang{15}/10.0 = \qty{329.8}{s}$, against
\qty{369.6}{s}; at \ang{90}: $2\times6371\times\sin\ang{45}/10.0 =
\qty{901.0}{s}$, against \qty{779.7}{s}. (c) The ray to \ang{30} turns in the
upper mantle, where the true speed is about 8 to \qty{10}{km.s^{-1}}, below the
uniform \qty{10}{km.s^{-1}}; the ray to \ang{90} spends most of its path in the
lower mantle, at 11 to \qty{13.7}{km.s^{-1}}, and bends down into it, so the
real Earth is faster.
\end{solution}

\begin{solution}{exo:b3:seismology-imaging:6}
(a) Squaring, $(2 - \xi)^4 = 16(1 - \xi/3)(1 - \xi)$. Expand:
$16 - 32\xi + 24\xi^2 - 8\xi^3 + \xi^4 = 16 - \tfrac{64}{3}\xi + \tfrac{16}{3}\xi^2$,
so $\xi^4 - 8\xi^3 + \tfrac{56}{3}\xi^2 - \tfrac{32}{3}\xi = 0$, that is
$\xi\,(3\xi^3 - 24\xi^2 + 56\xi - 32) = 0$, and
$3\xi^3 - 24\xi^2 + 56\xi - 32 = (\xi - 4)(3\xi^2 - 12\xi + 8)$. The root
$\xi = 0$ is not a wave. (b) The roots of the quadratic are $2 \pm 2/\sqrt3$,
that is 3.155 and 0.8453; only $\xi = 0.8453$ lies in $(0, 1)$ (and $\xi = 4$
does not), and it satisfies the unsquared equation, both sides being $4/3$.
So $c/\beta = \sqrt{0.8453} = 0.9194$. (c) $c = 0.9194\times3.46 =
\qty{3.181}{km.s^{-1}}$; over \qty{1000}{km}, $314.4 - 289.0 = \qty{25.3}{s}$
after the S wave.
\end{solution}

\begin{solution}{exo:b3:seismology-imaging:7}
$\mathsf G$ has the rows $(1, x_i)$ and $\mathsf C_d^{-1} = \operatorname{diag}
(100, 100, 25, 25)$ (\unit{s^{-2}}). The normal equations are
\[
  \begin{pmatrix} 250 & 23750 \\ 23750 & \num{2812500} \end{pmatrix}
  \begin{pmatrix} a \\ b \end{pmatrix}
  = \begin{pmatrix} 3665 \\ 418125 \end{pmatrix},
\]
with determinant \num{1.390625e8}. Hence $b = (250\times418125 -
23750\times3665)/\num{1.390625e8} = \qty{0.12575}{s.km^{-1}}$ and
$a = (2812500\times3665 - 23750\times418125)/\num{1.390625e8} =
\qty{2.713}{s}$. The covariance is $\mathsf N^{-1}$: $\sigma_b =
\sqrt{250/\num{1.390625e8}} = \qty{0.00134}{s.km^{-1}}$, $\sigma_a =
\sqrt{2812500/\num{1.390625e8}} = \qty{0.142}{s}$. The speed is
$1/b = \qty{7.95}{km.s^{-1}}$, with the error $\sigma_b/b^2 =
\qty{0.08}{km.s^{-1}}$: an uppermost-mantle speed. The residuals are $-0.001$,
$0.011$, $-0.076$ and $\qty{0.036}{s}$, and $\chi^2 = 0.19$ for $4 - 2 = 2$
degrees of freedom: smaller than expected, so the stated errors are, if
anything, generous; the fit is acceptable.
\end{solution}

\begin{solution}{exo:b3:seismology-imaging:8}
(a) $\mathsf G = (30\;\;40)$ is a single row: $\mathsf G^{\mathsf T}\mathsf G$ is
$2\times2$ of rank one, hence singular; any model with $30m_1 + 40m_2 = 0$ can
be added without changing the datum. (b) $\mathsf G\mathsf G^{\mathsf T} =
900 + 1600 = 2500$, so $\mathbf m = (30, 40)\times0.50/2500 =
(0.006, 0.008)\,\unit{s.km^{-1}}$. (c) $\mathsf R = \frac{1}{2500}
\begin{pmatrix} 900 & 1200 \\ 1200 & 1600 \end{pmatrix} =
\begin{pmatrix} 0.36 & 0.48 \\ 0.48 & 0.64 \end{pmatrix}$, of trace 1: one
datum resolves one combination of parameters. (d) $\mathsf R\,(0.01, 0) =
(0.0036, 0.0048)\,\unit{s.km^{-1}}$: the anomaly is spread along the ray, more
into the cell where the ray is longer.
\end{solution}

\begin{solution}{exo:b3:seismology-imaging:9}
(a) Linearised: $\delta t = -(\delta v/v)L/v = -0.03\times200/9.0 =
\qty{-0.667}{s}$; exactly, $L/(1.03v) - L/v = 200/9.27 - 200/9.0 =
\qty{-0.647}{s}$. (b) The path in the slab is $100/\cos\ang{40} =
\qty{130.5}{km}$: $-0.03\times130.5/9.0 = \qty{-0.435}{s}$ linearised and
$\qty{-0.422}{s}$ exactly. (c) The exact value is the linearised one divided by
$1.03$: an error of \qty{3}{\percent}, of the order of $\delta v/v$ itself, as
expected of a first-order theory.
\end{solution}

\begin{solution}{exo:b3:seismology-imaging:10}
(a) The overtone exists while its half-space part decays, $c < \beta_2$; at its
cut-off $c = \beta_2$, so $q_2 = 0$ and $\tan(kHq_1) = 0$, with $kHq_1$ in
$[n\pi, n\pi + \pi/2)$: $kHq_1 = n\pi$. (b) With $k = \omega/\beta_2$ and
$q_1 = \sqrt{\beta_2^2/\beta_1^2 - 1} = \sqrt{\beta_2^2 - \beta_1^2}/\beta_1$,
$\omega_nH\sqrt{\beta_2^2 - \beta_1^2}/(\beta_1\beta_2) = n\pi$, which is the
formula. (c) $\sqrt{4.5^2 - 3.5^2} = \sqrt8 = 2.828$; $\omega_1 =
\pi\times3.5\times4.5/(35\times2.828) = \qty{0.4998}{rad.s^{-1}}$, a period of
$2\pi/0.4998 = \qty{12.57}{s}$: at longer periods only the fundamental mode
exists. (d) For $n = 0$ the cut-off frequency is zero: the fundamental mode
exists at all periods. At short periods $kH \to \infty$ while $kHq_1$ stays
below $\pi/2$, so $q_1 \to 0$ and $c \to \beta_1$.
\end{solution}

\begin{solution}{exo:b3:seismology-imaging:11}
(a) $\mathsf G^{\mathsf T}(\mathsf G\mathsf G^{\mathsf T} + \varepsilon^2\mathsf I)
= (\mathsf G^{\mathsf T}\mathsf G + \varepsilon^2\mathsf I)\mathsf G^{\mathsf T}$;
multiply on the left by $(\mathsf G^{\mathsf T}\mathsf G + \varepsilon^2
\mathsf I)^{-1}$ and on the right by $(\mathsf G\mathsf G^{\mathsf T} +
\varepsilon^2\mathsf I)^{-1}$, both invertible for $\varepsilon > 0$.
(b) Here $\mathsf G\mathsf G^{\mathsf T} = \|\mathbf g\|^2 = 2500$, a number,
so $\mathsf R_\varepsilon = \mathsf G^{\mathsf T}\mathsf G/(2500 + \varepsilon^2)$,
of trace $2500/(2500 + \varepsilon^2)$. (c) The predicted datum is
$\mathbf g\cdot\mathbf m_\varepsilon = 0.50\times2500/(2500 + \varepsilon^2)$
and the residual $0.50\,\varepsilon^2/(2500 + \varepsilon^2)$. For
$\varepsilon^2 = 0$: trace 1, residual 0; for 2500: trace 0.5, residual
\qty{0.25}{s}; for \num{10000}: trace 0.2, residual \qty{0.40}{s}. More damping
gives a smaller model, a poorer resolution and a larger misfit.
\end{solution}

\begin{solution}{exo:b3:seismology-imaging:12}
(a) By exercise 5, $T = (2R/v)\sin(\Delta/2)$, so $\mathrm dT/\mathrm d\Delta
= (R/v)\cos(\Delta/2)$. A chord making the angle $\Delta$ at the centre is the
base of an isosceles triangle with apex angle $\Delta$, so it meets the radius
at the source at $i_0 = (\pi - \Delta)/2$, and $p = R\sin i_0/v =
(R/v)\cos(\Delta/2)$: the same. (b) The turning point is the middle of the
chord, at the distance $R\cos(\Delta_1/2)$ from the centre; equivalently
$r_1 = p_1v$. (c) Here $p(\Delta)/p_1 = \cos(\Delta/2)/\cos(\Delta_1/2)$, and
the formula gives $r_1 = R\exp\bigl(-\ln(1/\cos(\Delta_1/2))\bigr) =
R\cos(\Delta_1/2)$, as in (b); the speed is $r_1/p_1 = R\cos(\Delta_1/2)\big/
\bigl((R/v)\cos(\Delta_1/2)\bigr) = v$.
\end{solution}

\begin{solution}{pb:b3:seismology-imaging:1}
\textbf{1.} $i_0 = \arcsin 0.8 = \ang{53.13}$ and $p = 0.8/5.0 =
\qty{0.16}{s.km^{-1}}$.

\textbf{2.} The ray turns where $v = 1/p = \qty{6.25}{km.s^{-1}}$, that is at
$z = (6.25 - 5.0)/0.05 = \qty{25}{km}$.

\textbf{3.} Differentiating $\sin i = pv$ along the ray, $\cos i\,\mathrm di/
\mathrm ds = pg\,\mathrm dz/\mathrm ds = pg\cos i$, so $\mathrm di/\mathrm ds =
pg$: the direction turns at a constant rate, the curvature of a circle of radius
$1/(pg) = 1/(0.16\times0.05) = \qty{125}{km}$. At the turning point the ray is
horizontal at \qty{25}{km} depth and the centre lies vertically \qty{125}{km}
above it: at $z = \qty{-100}{km}$, that is $v_0/g$ above the surface, where the
extrapolated speed would vanish.

\textbf{4.} At the source the radius of the circle is perpendicular to the ray,
so it makes the angle $i_0$ with the horizontal, and the centre lies
$125\cos i_0$ away horizontally; by symmetry
$X = 2\times125\times\cos i_0 = 2\times125\times0.6 = \qty{150}{km}$, or
$X = 2\sqrt{1 - p^2v_0^2}/(pg)$.

\textbf{5.} $\mathrm dT = \mathrm ds/v$ with $\mathrm ds = \mathrm di/(pg)$ and
$v = \sin i/p$: $\mathrm dT = \mathrm di/(g\sin i)$. Then
\[
  T = \frac2g\int_{i_0}^{\pi/2}\frac{\mathrm di}{\sin i} = \frac2g\ln\cot\frac{i_0}2 =
  \frac2g\operatorname{arccosh}\frac1{pv_0} .
\]
Here $\cot(\ang{26.57}) = 2$, so
$T = 40\ln2 = \qty{27.73}{s}$.

\textbf{6.} For $p = 0.15$: $pv_0 = 0.75$, $\sqrt{1 - 0.5625} = 0.6614$,
$X = 2\times0.6614/0.0075 = \qty{176.4}{km}$ and $T = 40\operatorname{arccosh}
(1.3333) = 40\times0.7954 = \qty{31.81}{s}$. Then $\Delta T/\Delta X =
4.08/26.4 = \qty{0.155}{s.km^{-1}}$, between 0.15 and 0.16, as
$\mathrm dT/\mathrm dX = p$ requires.

\textbf{7.} $\mu_1 = 2700\times3500^2 = \qty{33.1}{GPa}$ and
$\mu_2 = 3300\times4500^2 = \qty{66.8}{GPa}$.

\textbf{8.} $k = 2\pi/(cT) = 2\pi/(3.99\times25) = \qty{0.06299}{rad.km^{-1}}$;
$q_1 = \sqrt{(3.99/3.5)^2 - 1} = 0.5474$; $q_2 = \sqrt{1 - (3.99/4.5)^2} = 0.4624$.

\textbf{9.} $\tan(kHq_1) = 66.825\times0.4624/(33.075\times0.5474) = 1.707$; the
fundamental mode takes $kHq_1 = \arctan1.707 = 1.041$, so
$H = 1.041/(0.06299\times0.5474) = \qty{30.2}{km}$.

\textbf{10.} At \qty{40}{s}: $k = 2\pi/170 = 0.03696$, $q_1 = 0.6888$,
$q_2 = 0.3287$, $\tan(kHq_1) = 0.964$, $kHq_1 = 0.7671$ and
$H = 0.7671/(0.03696\times0.6888) = \qty{30.1}{km}$. The two agree within the
precision of the phase velocities: the crust is about \qty{30}{km} thick.

\textbf{11.} $\omega = 0.2513$ and $\qty{0.1571}{rad.s^{-1}}$, $k = 0.06299$ and
$\qty{0.03696}{rad.km^{-1}}$: $U \approx 0.0942/0.02603 = \qty{3.62}{km.s^{-1}}$,
less than both phase velocities, as it must be when $c$ increases with the
period.

\textbf{12.} $\lambda = 3.99\times25 = \qty{99.8}{km}$ and $4.25\times40 =
\qty{170}{km}$. The fundamental mode samples down to a fraction of its
wavelength: the 25-s wave lies mostly in the crust, the 40-s wave reaches well
into the mantle, so its speed is closer to $\beta_2$. Two periods fixing the
same $H$ is a check that the two-layer model is adequate.

\textbf{13.} With the columns for cells 1 to 4 and the rows A to F, in units
of \qty{1000}{km} ($\sqrt2\times100/1000 = 0.1414$):
\[
  \mathsf G = \begin{pmatrix}
  0.1 & 0.1 & 0 & 0 \\ 0 & 0 & 0.1 & 0.1 \\ 0.1 & 0 & 0.1 & 0 \\
  0 & 0.1 & 0 & 0.1 \\ 0.1414 & 0 & 0 & 0.1414 \\ 0 & 0.1414 & 0.1414 & 0
  \end{pmatrix}.
\]

\textbf{14.} $\mathsf C_d^{-1} = \operatorname{diag}(400, 400, 400, 400, 100,
100)$. A row or column ray contributes $0.1^2\times400 = 4$, a diagonal ray
$0.1414^2\times100 = 2$. Cell 1 is crossed by A, C (4 each) and E (2): 10. Cells
1 and 2 share ray A: 4; cells 1 and 3 share C: 4; cells 1 and 4 share E: 2; and
so on by symmetry, which gives the four rows stated.

\textbf{15.} A row or column ray contributes $0.1\times400\,\delta t = 40\,
\delta t$ to each of its cells, a diagonal ray $0.1414\times100\,\delta t =
14.14\,\delta t$. Cell 1: $40(0.16 - 0.36) + 14.14\times(-0.15) = -10.12$; cell
2: $40(0.16 + 0.11) + 14.14\times(-0.27) = 6.98$; cell 3: $40(-0.48 - 0.36) +
14.14\times(-0.27) = -37.42$; cell 4: $40(-0.48 + 0.11) + 14.14\times(-0.15) =
-16.92$.

\textbf{16.} For $(1,1,1,1)$, each row of $\mathsf N$ sums to 20. For
$(1,-1,-1,1)$, row 1 gives $10 - 4 - 4 + 2 = 4$ and the others the same
pattern. For $(1,1,-1,-1)$, $10 + 4 - 4 - 2 = 8$, and for $(1,-1,1,-1)$,
$10 - 4 + 4 - 2 = 8$. The vectors are orthogonal with norm 2, so
$\mathsf N^{-1} = \sum_k \mathbf u_k\mathbf u_k^{\mathsf T}/(4\lambda_k)$; its
$(1,1)$ entry is $\tfrac14(\tfrac1{20} + \tfrac14 + \tfrac18 + \tfrac18) =
\tfrac{11}{80}$, its $(1,2)$ entry $\tfrac14(\tfrac1{20} - \tfrac14 + \tfrac18
- \tfrac18) = -\tfrac4{80}$, its $(1,4)$ entry $\tfrac14(\tfrac1{20} +
\tfrac14 - \tfrac18 - \tfrac18) = \tfrac1{80}$, and so on.

\textbf{17.} $\hat{\mathbf m} = \mathsf N^{-1}\mathbf b$:
$\hat m_1 = (11\times(-10.121) - 4\times6.982 - 4\times(-37.418) - 16.921)/80
= -0.081$; likewise $\hat m_2 = 1.844$, $\hat m_3 = -3.706$ and
$\hat m_4 = \qty{-0.931}{ms.km^{-1}}$.

\textbf{18.} $\delta v/v_0 = -v_0\,\delta s = -8.0\times10^{-3}\,\hat m$, in per
cent $-0.8\,\hat m$: cell 1 $+0.07$, cell 2 $-1.48$, cell 3 $+2.96$, cell 4
$\qty{+0.75}{\percent}$.

\textbf{19.} Every diagonal entry of $\mathsf N^{-1}$ is $11/80$: each cell has
the standard error $\sqrt{0.1375} = \qty{0.371}{ms.km^{-1}}$, that is
$\qty{0.30}{\percent}$. The correlation of cells 1 and 2 is $-4/11 = -0.36$:
they share ray A, and a slowness put in one can be traded against the other.
Cells 1 and 4, which share only the less precise diagonal E, are weakly and
positively correlated, $+1/11 = +0.09$.

\textbf{20.} The predicted residuals are $0.176$, $-0.464$, $-0.379$, $0.091$,
$-0.143$ and $\qty{-0.263}{s}$; the residuals after the fit are $-0.016$,
$-0.016$, $0.019$, $0.019$, $-0.007$ and $\qty{-0.007}{s}$, and
$\chi^2 = 0.50$ for $6 - 4 = 2$ degrees of freedom: consistent with the stated
errors, if anything on the small side.

\textbf{21.} The eigenvectors are those of $\mathsf N$ and the eigenvalues
$\lambda/(\lambda + 4)$: $20/24 = 5/6$, $4/8 = 1/2$, $8/12 = 2/3$ twice. As in
question 16, $R_{11} = \tfrac14(\tfrac56 + \tfrac12 + \tfrac23 + \tfrac23) =
\tfrac23$; $R_{12} = \tfrac14(\tfrac56 - \tfrac12 + \tfrac23 - \tfrac23) =
\tfrac1{12}$; $R_{13} = \tfrac14(\tfrac56 - \tfrac12 - \tfrac23 + \tfrac23) =
\tfrac1{12}$; $R_{14} = \tfrac14(\tfrac56 + \tfrac12 - \tfrac23 - \tfrac23) = 0$.

\textbf{22.} Project $\mathbf b$ on the eigenvectors: $\mathbf b\cdot\mathbf u_k/4
= -14.37$, $0.849$, $12.80$ and $-9.40$ for the eigenvalues 20, 4, 8 and 8.
The damped estimate divides them by $\lambda + 4$ instead of $\lambda$; for
cell 3 (components $1, -1, -1, 1$),
$m_3 = -14.37/24 - 0.849/8 - 12.80/12 - 9.40/12 = \qty{-2.555}{ms.km^{-1}}$,
a speed anomaly of $\qty{+2.04}{\percent}$. Damping pulls every component
towards zero, most of all the poorly constrained ones.

\textbf{23.} The undamped model, which fits the data to within their errors
($\chi^2 = 0.50$ for 2 degrees of freedom), whereas the damped one misfits them
badly ($15.8$). Damping is needed only when $\mathsf N$ is singular or nearly
so; here its smallest eigenvalue, 4, is not small, the estimate is stable and
its errors are moderate.

\textbf{24.} Row 3 is $(1/12, 0, 2/3, 1/12)$: the damped estimate of cell 3 is
two thirds of its true value plus a twelfth of cells 1 and 4, its neighbours
along ray C and ray B, and nothing of cell 2.

\textbf{25.} Cell 3, the lower cell on the left, is fast by about
\qty{3}{\percent}; cell 2 is slow. A cold, dense slab of oceanic lithosphere
sinking beneath the continent is faster than the surrounding mantle, so the
fast cell, deep and on the ocean side of K, is the slab; the slow cell above
it on the other side lies where the mantle wedge is hot.

\textbf{26.} From questions 18 and 19, the P-velocity anomaly of the slab cell is
\textbf{\qty[text-series-to-math=true, uncertainty-mode=separate]{3.0 \pm 0.3}{\percent}}.
\end{solution}
